\chapter{Installation and Usage Notes}

\section{Python Environment}

The recommended setup is:

\begin{lstlisting}[style=pythonstyle]
python -m venv venv
venv\Scripts\activate
pip install -r requirements.txt
\end{lstlisting}

If PyTorch CPU wheels must be installed separately on Windows, the following command can be used before installing the remaining requirements:

\begin{lstlisting}[style=pythonstyle]
pip install torch torchvision --index-url https://download.pytorch.org/whl/cpu
\end{lstlisting}

\section{Typical Workflow}

\begin{enumerate}
  \item Collect employee images:
\begin{lstlisting}[style=pythonstyle]
python src/data/collect_faces.py
\end{lstlisting}
  \item Augment images:
\begin{lstlisting}[style=pythonstyle]
python src/data/augment.py
\end{lstlisting}
  \item Build the face database:
\begin{lstlisting}[style=pythonstyle]
python src/recognition/build_database.py
\end{lstlisting}
  \item Run the backend:
\begin{lstlisting}[style=pythonstyle]
uvicorn src.api.main:app --reload --port 8000
\end{lstlisting}
  \item Run the frontend:
\begin{lstlisting}[style=pythonstyle]
python -m http.server 5500 --directory frontend
\end{lstlisting}
  \item Open the dashboard at:
\begin{lstlisting}[style=pythonstyle]
http://localhost:5500
\end{lstlisting}
\end{enumerate}

\section{Main API Examples}

\begin{lstlisting}[style=pythonstyle]
GET  /health
POST /api/employees/register
GET  /api/employees
POST /api/recognize
GET  /api/attendance/report
POST /api/attendance/manual
GET  /api/attendance/logs
GET  /api/settings
POST /api/settings
\end{lstlisting}

\section{Suggested Evaluation Protocol}

For future formal evaluation, the following protocol is recommended:

\begin{enumerate}
  \item Register each employee with 5 to 10 high-quality enrollment images.
  \item Capture at least 20 validation images per employee under realistic conditions.
  \item Include impostor attempts from people not enrolled in the database.
  \item Compute genuine and impostor similarity distributions.
  \item Sweep the recognition threshold from 0.2 to 0.9.
  \item Report true acceptance, false acceptance, false rejection, and equal error rate.
  \item Repeat the test under different lighting conditions and camera distances.
\end{enumerate}

\chapter{Selected Code Excerpts}

\section{Recognition Decision Logic}

The recognizer normalizes embeddings, searches the database, and compares the best similarity with a runtime threshold. The simplified logic is:

\begin{lstlisting}[style=pythonstyle,language=Python]
emb = normalize(input_embedding)
sims, idxs = search(emb, k=min(5, len(labels)))
best_sim = float(sims[0])
threshold = get_runtime_settings()["face_threshold"]

if best_sim < threshold:
    return {
        "identity": "unknown",
        "confidence": best_sim,
        "status": "unknown"
    }

top_labels = [labels[i] for i in idxs if i >= 0]
identity = Counter(top_labels).most_common(1)[0][0]
return {
    "identity": identity,
    "confidence": best_sim,
    "status": "recognized"
}
\end{lstlisting}

\section{Attendance Event Decision}

The automatic attendance service uses the first event of the day as check-in and a later event as check-out, while blocking repeated recognition during the cooldown period.

\begin{lstlisting}[style=pythonstyle,language=Python]
recent = query_recent_log(employee_id, cooldown_minutes)
if recent:
    return {"logged": False, "reason": "cooldown"}

has_check_in = query_today_check_in(employee_id)
event_type = "check_out" if has_check_in else "check_in"

create_attendance_log(
    employee_id=employee_id,
    event_type=event_type,
    confidence=confidence,
    camera_id=camera_id,
    is_real_face=is_real,
    liveness_score=liveness_score
)
\end{lstlisting}
